In this section, we present \lang's formal semantics. We first discuss the terms
in the language (\Sec{sec:terms}), then the types (\Sec{sec:types}) and
\emph{regions} (\Sec{sec:regions}), and our environments and the
mechanics of typechecking (\Sec{sec:typechecking}). Finally, we move on to
discussion of our metatheory (\Sec{sec:metatheory}) and tested semantics
(\Sec{sec:tested-semantics}).

\subsection{The Syntax of \lang}
\label{sec:terms}

\begin{figure}
  \figuresize
  \oxnamegrammar

  \begin{minipage}{.43\textwidth}
    \begin{bnf}
      \oxpathgrammar \\
      \oxplacegrammar \\
      \oxplaceexprgrammar \\
      \oxplaceexprctxgrammar \\[1em]

      \oxprovgrammar \\
      \oxmutagrammar \\[1em]
    \end{bnf}
  \end{minipage}
  \begin{minipage}{.56\textwidth}
    \begin{bnf}
      \oxprimitivegrammar \\
      \oxexpressiongrammar \\
    \end{bnf}
  \end{minipage}

  \vspace{-1em}
  \caption{Term Syntax of \lang}
  \label{fig:term-syntax}
  \vspace{-1em}
\end{figure}

\Fig{fig:term-syntax} presents the syntax of \lang terms, in four broad
groupings: (1)~metavariables for the various kinds of names that exist,
(2)~places, which act as names for abstract memory locations, 
(3)~annotations for references, and (4)~the actual terms of the language. The
first group is fairly conventional, but we'll discuss special names like
frame variables as they come up.

\paragraph{Places and Place Expressions} As we saw in \Sec{sec:borrowing},
places $\oxplace$ and place expressions $\oxplaceexpr$ are names for paths from
a particular variable to a particular part of the object stored there, whether
that be a projection of a tuple, or a field of a struct (where a struct is
really precisely just a named tuple or record type). One might think of place
expressions as a sort of syntactic generalization of variables. They are
analogous to what are called \texttt{lvalues} in C. Places $\oxplace$ are a
subset of place expressions that do not include dereferences. They can
intuitively be thought of as an abstract name of a memory location since when
bound, they will always correspond to one particular value on the stack. Place
expression contexts $\oxplaceexprctx$ are used in various parts of the formalism
to decompose place expressions $\oxplaceexpr$ into an innermost dereferenced
place, $\oxderef{\oxplace}$, and an outer context $\oxplaceexprctx$.

\paragraph{Annotations for References} In \lang, we have two annotations that we
provide for every borrowing expression. First, we annotate references with
ownership qualifiers $\oxmuta$, indicating whether the reference is shared
($\oximm$) or unique ($\oxmut$). We use these rather than their equivalents in
Rust (no annotation and $\rusti{mut}$ respectively) because the terms more
accurately reflect the semantic focus on \emph{aliasing}, rather than
\emph{mutation}. Indeed, in Rust, a value of the type \rusti{&&mut u32}
\emph{cannot} be mutated (because we have a shared reference to a unique
reference), and a value of the type \rusti{&Cell<u32>}\footnote{\rusti{Cell<T>}
  is a Rust standard library type that provides a ``mutable memory location''
  that allows mutation in its API.} \emph{can} be mutated through the method
\rusti{Cell::set}.

Second, we annotate references with \emph{regions}. Regions $\oxprov$ have two
forms: abstract regions $\oxabsprov$ (pronounced \emph{var-rho}) and concrete
regions $\oxcprov$. Abstract regions correspond to lifetime variables
\rusti{'a}, \rusti{'b}, etc.\ in Rust, and are used polymorphically in function
types to indicate that the function is agnostic to the particular regions of
reference-type parameters. Concrete regions, by contrast, carry concrete
information in the environment where they correspond to a set of loans. A loan
$\oxloanpkg{\oxmuta}{\oxplaceexpr}$ indicates a possible origin
($\oxplaceexpr$), qualified by whether the loan is unique or shared ($\oxmuta$).
Intuitively, each loan tells us a single possible origin for a reference, while
a concrete region maps to \emph{all} possible origins of a reference. As we will
see in \Sec{sec:regions}, regions are essential to enabling our type system to
guarantee the correct use of unique and shared references.

\paragraph{Expressions} Expressions $\oxexpr$ in \lang are numerous, but largely
standard. For example, constants $\oxprim$ consist of the unit value $\oxunit$,
unsigned 32-bit integers $\oxnum$, and boolean values $\oxtrue$ and $\oxfalse$.
The most interesting expressions in \lang are the ones we've already seen by
example: place expression usage (written simply $\oxplaceexpr$) and borrowing
(with several forms that we explain shortly). The former may be thought of as
variables that behave linearly for non-copyable data (removing the place from
the environment after use), and traditionally for copyable data. (As a first
approximation, one can think of all data that is not a unique pointer as safely
copyable.)

There are three borrowing forms, and all work in fundamentally the same way:
they are each used as introduction forms for references. The simplest case,
written $\oxref{\oxcprov}{\oxmuta}{\oxplaceexpr}$, introduces an
$\oxmuta$-reference (with region $\oxcprov$) to the location that the place
expression $\oxplaceexpr$ evaluates to. The next form borrows from
$\oxindex{\oxplaceexpr}{\oxexpr}$ instead of simply $\oxplaceexpr$, and is used
to borrow an element out of an array or slice $\oxplaceexpr$ at the index given
by $\oxexpr$. The final form borrows from
$\oxslice{\oxplaceexpr}{\oxexpr_1}{\oxexpr_2}$, and is used to borrow a slice of
$\oxplaceexpr$ using the range given by $\oxexpr_1$ and $\oxexpr_2$. A
\emph{slice} is Rust terminology for a dynamically-sized subsection of an array.

In these last two cases, one might wonder ``why are indexing and slicing not
places themselves?'' The answer comes in two parts: (1)~indexing and slicing
take arbitrary expressions, while places are entirely static, and (2)~unlike
tuple projections which have a fine-grained notion of ownership, indexing and
slicing affect the ownership of the array or slice overall. This second part
means that while you can create two unique references to different projections
of the same tuple, you cannot create two unique references to different indices
of an array.

The remainder of our expressions are standard or discussed already. These
include sequencing, assignment, and creation of tuples and arrays. Our closure
syntax follows the syntax of Rust, and thus uses vertical bars to denote the
closure's parameters. As in Rust, closures are not polymorphic; only global
functions (shown in \Fig{fig:envs}) may be polymorphic and specify where-bounds
on regions.\footnote{In Rust, where-bounds in functions are used to constrain one
  lifetime to outlive another, meaning that a reference with the larger lifetime
  must be valid at least as long as a reference with the shorter lifetime.} We
use function application when applying closures as well as global functions.
Hence, function application additionally includes polymorphic instantiation
written using Rust's turbofish syntax (\rusti{::<>}). An $\oxabort{\oxstring}$
indicates irrecoverable failure; it terminates the program with the given string
as a diagnostic message. Finally, \lang includes tagged sums, which are
introduced using the \texttt{Left} and \texttt{Right} forms and eliminated using
\texttt{match}.\footnote{Rust, of course, supports more general n-ary tagged sums
  with user-definable tags (calling the whole system enumerations), but binary
  sums suffice to get at the essence of Rust without requiring a complicated
  formalization for pattern matching}


\subsection{Types in \lang}
\label{sec:types}

\begin{figure}
  \figuresize
  \begin{bnf}
    \oxkindgrammar \\
    \oxbasetypegrammar \\
  \end{bnf}
  \begin{bnf}
    \oxsitypegrammar \\
  \end{bnf}
  \begin{bnf}
    \oxxitypegrammar \\
    \oxsdtypegrammar \\
    \oxsxtypegrammar \\
    \oxtypegrammar \\
  \end{bnf}

  \vspace{-1em}
  \caption{Type Syntax of \lang}
  \label{fig:type-syntax}
  \vspace{-1em}
\end{figure}

In \lang, we have five distinct categories of types (based on two features we
need to distinguish: sized vs. unsized, and initialized vs. dead), and a kind
system to track the three kinds of polymorphism in the language. While these
distinctions may seem complex, they greatly simplify the well-formedness
conditions required on types during typechecking. The grammars are all present
in \Fig{fig:type-syntax}, and are explained in detail in the rest of the
section.

\paragraph{Sized and Unsized Types} We need to distinguish between types based
on sizedness, which is a direct consequence of Rust itself. All bindings in Rust
(and in \lang) must be able to fit on the stack which requires that they have a
statically-known size. In Rust, this is dealt with using a special
automatically-derived \emph{marker trait} called \rusti{Sized} which serves as a
tag during typechecking to indicate that a type has a statically-knowable size.
For pragmatic reasons (since one typically works with sized types), Rust decided
on using \rusti{Sized?} to indicate that a type is ``possibly unsized'' (and
thus could only be part of a type for a let binding if it is behind a
reference). In \lang, we have a comparable syntactic distinction between
\emph{sized} types $\oxsitype$ and \emph{maybe unsized} types $\oxxitype$. Sized
types characterize all the types with statically-known sizes and maybe unsized
types $\oxxitype$ include all such types \emph{and} the slice type
$\oxtslice{\oxsitype}$ which corresponds to a dynamically-sized portion of an
array.

\paragraph{Initialized and Dead Types} We also need to distinguish between types
based on initialization, which we use to model the 'use-once' linearity of
variables referring to non-copyable data. To that end, we introduce two
categories. First, \emph{dead} types $\oxsdtype$ which is either a sized and
initialized type with a dagger on it (indicating that it is dead) or a product
of dead types. These correspond to totally moved types. Second, we have
\emph{maybe dead} types $\oxsxtype$ which can be either initialized, dead, or a
product of maybe dead types, corresponding to types where some of their
components have been moved. Though not supported directly in our formalism,
these dead and maybe dead types also can be used directly to support
uninitialized and partially-initialized variable bindings.

\paragraph{Kinds and Polymorphism} \lang has three kinds $\oxkind$: the kind of
ordinary types $\oxktype$, the kind of regions $\oxkprov$, and the kind of
\emph{frame typings} $\oxkenv$. (Frame typings are relevant for closures, as
we'll see below.) We abstract over variables of each kind in \lang and, to aid
the reader, we have separate syntax for each: $\oxtvar$, $\oxabsprov$, and
$\oxenvvar$, respectively. For simplicity, \lang restricts type variables
$\oxtvar$ to being instantiated only with sized and initialized types, but this
limitation could be addressed by enriching kinds further with a unique kind for
each sort of type.

\paragraph{The Types Themselves} The majority of types in \lang are sized \&
initialized types, including base types $\oxbasetype$, type variables $\oxtvar$,
tuples $\oxtprod{\oxsitype_1 \oxdotsc \oxsitype_n}$, arrays of length $n$
$\oxtarr{\oxsitype}{\oxnum}$, binary sums $\oxtsum{\oxsitype_1}{\oxsitype_2}$,
references $\oxtref{\oxprov}{\oxmuta}{\oxxitype}$, and function types.
% $\oxtfunext{ \overline{\oxenvvar} \oxcomma \overline{\oxabsprov} \oxcomma \overline{\oxtvar} }{\oxsitype_1 \oxdotsc \oxsitype_n}{\oxsitype_r\;}{\oxenv}{ \;\overline{\oxabsprov_1: \oxabsprov_2} }$.
With the exception of references, any types that occur within these types are
themselves required to be both sized and initialized. For reference types
$\oxtref{\oxprov}{\oxmuta}{\oxxitype}$, we include both the region $\oxprov$ and
ownership qualifier $\oxmuta$ in the type which allow us to understand
statically both a reference's origin as well as its aliasing requirements. We
allow potentially unsized types under references since the reference itself will
always have a fixed size regardless of what it points to (e.g. 64-bit on a
64-bit machine).

Function types have three notable features. First, each function type can
possibly include a frame expression $\oxenv$ (syntax in \Fig{fig:envs}) over the
arrow indicating what bindings, if any, were caught up in the closure
environment (when nothing is captured, we put nothing over the arrow). Next,
functions are polymorphic in type and region variables, as well as in frame
variables $\oxenvvar$ to enable the use of higher-order functions. Finally,
functions can relate types with abstract regions using outlives bounds:
\texttt{where} $\oxabsprov_1: \oxabsprov_2$ means $\oxabsprov_1$ outlives
$\oxabsprov_2$. These \texttt{where} bounds come directly from Rust, and are
useful in making functions that, e.g., reborrow from one of several
reference-typed parameters.

\begin{figure}
  \figuresize
  \begin{bnf}
    \oxglobalctxgrammar \\
    \oxglobalentrygrammar \\
  \end{bnf}
  \begin{bnf}
    \oxtvarctxgrammar \\
    \oxkontctxgrammar \\
    \oxvarctxgrammar \\
    \oxframegrammar \\
    \oxenvgrammar \\
  \end{bnf}

  \vspace{-1em}
  \caption{Environments in \lang}
  \label{fig:envs}
  \vspace{-1em}
\end{figure}

\subsection{Environments for Typechecking}
\label{sec:envs}

With the syntax of terms and types in hand, we can look more closely at some
example Oxide programs to understand the environments we'll be using for
typechecking. We'll start with a simple example using reborrowing, much like our
last example in \Sec{sec:main}.

\begin{rustbk}
  struct Obj(u32);
  letrgn<'y, 'z> { // 'y -> {}, 'z -> {}
    let mut x = Obj(5); // x : Obj
    let y = &'y uniq x /* 'y -> { $\oxloanpkg{\oxmut}{x}$ } */;
    let z = &'z uniq *y /* 'z -> { $\oxloanpkg{\oxmut}{x}$, $\oxloanpkg{\oxmut}{\oxderef{\! y}}$ } */;
  }
\end{rustbk}

Here, we create an object named \rusti{x} and in the comment, we see how our
stack typing (written $\oxvarctx$) will record the new binding and its type.
Then, on line 4, we produce a unique reference to \rusti{x}. In the comment, we
see the metadata produced by this borrow associating the region \rusti{'y} with
the set of loans $\oxset{\oxloanpkg{\oxmut}{x}}$. This metadata means that a
reference with the region \rusti{'y} must necessarily point to \rusti{x}. Note
that this metadata is produced immediately after the borrow expression, and
doesn't depend on the binding \rusti{y} being introduced for it. Then, on line
5, we reborrow from \rusti{y} as \rusti{z}, and again, we can see the
corresponding metadata produced associating the region \rusti{'z} with the set
of loans
$\oxset{ \oxloanpkg{\oxmut}{x} \oxcomma \oxloanpkg{\oxmut}{\oxderef{\! y}} }$.
Note that borrowing $\oxderef{y}$ here means reborrowing from \rusti{y}. This
tells us two things: (1)~that a reference with the region \rusti{'z} points to
\rusti{x}, and (2)~that a reference with the region \rusti{'z} was created by
reborrowing from \rusti{y}. That latter means that \rusti{y} ought to be
rendered unusable as long as our reference \rusti{z} (with region \rusti{'z})
exists. The two pieces of information we've seen here, in-scope bindings with
their types and borrowing metadata, are precisely what's necessary for us to
track in our stack typing $\oxvarctx$, and each entry follows the syntax seen
here. In this next example, we'll see a bit more complexity by defining and
using a closure.

\begin{rustbk}
  let x = Obj(5); // x : Obj
  let y = Obj(9); // y : Obj
  letrgn<'z> { // z -> {}
    let f = |obj: &'z uniq Obj| -> Obj { Obj((*obj).0 + y.0) }; // $\oxvarctxentry{y}{\oxtuninit{Obj}}$
    let z = &'z uniq x /* 'z -> { $\oxloanpkg{\oxmut}{x}$ } */;
    f(z) // $\oxvarctxentry{z}{\oxtuninit{(\oxtref{'z}{\oxmut}{Obj})}}$
  }
\end{rustbk}

We create two objects named \rusti{x} and \rusti{y} respectively. Then, on lines
4-6, we define a closure named \rusti{f} that moves \rusti{y} from the context.
This movement is described by the annotation on line 6 that shows the new entry
for \rusti{y} in our stack typing $\oxvarctx$. This new entry differs in that
the type associated with \rusti{y} is marked with a dagger indicating that it is
now dead. On line 7, we create a reference \rusti{z} that we then pass on line 8
as an argument for \rusti{f}. On line 7, we also see the annotation for the
region indicating that a reference with region \rusti{'z} must point to
\rusti{x}, much like we saw in the previous example. In the comment on line 8,
we see that the use of \rusti{z} \emph{moved} it into the function call as well,
thus we mark its whole type with the dagger.

\begin{figure}
  \figuresize
  \begin{flushleft}
    \oxmusafetyform where $\oxmusafetyshape$ means
    $\oxmusafetyinner{\oxtvarctx}{\oxkontctx}{\oxvarctx}{\oxmuta}{\oxemptyctx}{\oxplaceexpr}{ \oxset{\overline{\oxloanpkg{\oxmuta}{\oxplaceexpr^\prime}}} }$.
  \end{flushleft}

  \begin{mathpar}
    \OSafePlace \and \ODeref \and \ODerefAbs
  \end{mathpar}

  \vspace{-1em}
  \caption{Ownership Safety in \lang}
  \label{fig:safety}
  \vspace{-1em}
\end{figure}

One question that might arise here is ``given we marked it dead, how are we
still allowed to use \rusti{y} \emph{inside} the closure?'' This is an important
point that leads to a key aspect of the structure of our stack typing
$\oxvarctx$. The low-level nature of \lang means we need to actually statically
record the frame-based structure of the stack at runtime. So, when we construct
this closure, the moved objects (in this case, solely \rusti{y}) are recorded in
a new frame $\oxframe$ that is tracked in the type of the closure,
$\oxtfun{}{\oxtref{'z}{\;\oxmut\;}{\texttt{Obj}}}{\texttt{Obj}}{\oxframe}$. This
frame provides the type information necessary to typecheck the body in the
future, and the closure at runtime has a corresponding stack frame as part of
its value form. When we typecheck the body of the closure, they are appended to
the current stack typing with $\oxframesep$. So, our stack typing is really a
collection of frames $\oxframe$ separated by $\oxframesep$ and our frames
$\oxframe$ contain both in-scope bindings with their types and in-scope regions
with their associated loan sets. As the language of stacks and frames may imply,
both stack typings $\oxvarctx$ and frames $\oxframe$ are \emph{ordered} which
makes it easier to state invariants like outlives (which we will see in
\Sec{sec:typechecking}) and resolves many of the typical issues that arise
formally with variable binding.

This is all formally-defined in \Fig{fig:envs} which features a grammar for all
of the environments used in the static semantics. As suggested by the grammars,
there are two other environments that are non-conventional. First, we have a
global environment $\oxglobalctx$ which consists of a series of top-level
function definitions. We define the well-formedness of these function
definitions by saying that their bodies must be well-typed assuming that all
other functions (including itself) are well-typed at their annotated type, which
enables a simple treatment of mutually recursive function definitions. Second,
we have a temporary typing $\oxkontctx$ which consists of a sequence of types
for parts of objects that will be in the continuation of the term being
typechecked. The need for this is subtle, but we will explain it using the
following example:

\begin{rustbk}
  letrgn<'a, 'b> {
    let x = Obj(5);
    let y = Obj(9);
    let tup = (&'a uniq x, (); &'b uniq x);
  }
\end{rustbk}

In this example, we create two objects \rusti{x} and \rusti{y} and then attempt
to create a pair named \rusti{tup} that consists of two unique references.
Exactly as written, we first uniquely borrow \rusti{x} with region \rusti{'a}
and then in the second component, we sequence a unit value with a unique borrow
of \rusti{x} with region \rusti{'b}. Of course, the very idea of having a
product of unique references to the same data sounds like a contradiction and so
we would hope to reject this program! However, to capture the expressivity of
Rust's borrow-checking in \lang, we also have to clear loan sets associated with
unused regions at sequencing points in programs. This leads to a dilemma: by
reading the program, we know that the first reference with region \rusti{'a} is
still used when we go to define the second one, but the naive definition of use
would correspond to ``is there currently a reference with that region in the
stack typing?'' The continuation typing comes in to resolve this. We don't have
a name (yet) for the already-typechecked portions of a product, but we can
record their types in the continuation typing to record that they are around
since they may then be let-bound and used further in the program. We'll see
formally how this interaction happens during typechecking in the
\oxname{T-Tuple} rule in \Sec{sec:typechecking}.

Now, we can look at the shape of our typing judgment, which we will return to
define in \Sec{sec:typechecking}. The shape of our typing judgement is
$\oxtypjudgeshape$. This is read: with global environment $\oxglobalctx$, type
environment $\oxtvarctx$, temporary typing $\oxkontctx$, and stack typing
$\oxvarctx$, $\oxexpr$ is a well-typed expression of type $\oxtype$ with an
updated stack typing $\oxvarctx^\prime$ for use in typechecking the continuation
of $\oxexpr$.


\subsection{Region-Based Alias Management}
\label{sec:regions}

As discussed in \Sec{sec:terms}, borrow expressions in \lang all have
\emph{region} annotations which are used to associate references statically with
information about their possible referents. This information is essential to our
formulation of borrow checking since we leverage it to determine if new
references would be safe to create. This is done formally with a judgement
called \emph{ownership safety} which has the form $\oxmusafetyshape$. We can
read it as saying ``in the environments $\oxtvarctx$ and $\oxvarctx$, it is safe
to use the place expression $\oxplaceexpr$ $\oxmuta$-ly and producing all of the
loans in $\oxset{\overline{\oxloanpkg{\oxmuta}{\oxplaceexpr^\prime}}}$ (called
the borrow chain).,'' where $\oxmuta$ is either $\oxmut$ or $\oximm$. That is,
if we have a derivation where $\oxmuta$ is $\oxmut$, we know that we can use the
place expression $\oxplaceexpr$ uniquely because we have a proof that there are
no live loans against the section(s) of memory that $\oxplaceexpr$ represents.
Further, when we have a derivation where $\oxmuta$ is $\oximm$, we know that we
can use the place expression $\oxplaceexpr$ sharedly because we have a proof
that there are no live \emph{unique} loans against the section(s) of memory that
$\oxplaceexpr$ represents. The produced borrow chain is used to create the
borrowing metadata to store in the environment when the ownership safety check
is guarding a borrow expression, and in the simplest case, is just precisely
$\oxplaceexpr$.

Since it is precisely this ownership safety judgment that captures the essence
of Rust's ownership semantics, we understand Rust's borrow checking system as
ultimately being a system for statically building a proof that data in memory is
either \emph{uniquely owned} (and thus able to allow unguarded mutation) or
\emph{collectively shared}, but not both. To do so, intuitively, ownership
safety looks at all of the concrete regions in $\oxvarctx$, and ensures that all
the loans they map to are not in conflict with the place expression
$\oxplaceexpr$ we are attempting to use. For a $\oxmut$ borrow, a conflict
occurs if \emph{any} loan maps to an overlapping place, but for a $\oximm$
borrow, a conflict occurs only when a $\oxmut$ loan maps to an overlapping
place. Since places are abstract memory locations, we can consider two places as
overlapping if they are equal or one is a prefix of the other (meaning that the
the longer place corresponds to a piece of the larger object in memory that the
shorter place refers to).

Unfortunately, reborrowing complicates matters. To support reborrowing,
ownership safety uses an expanded inner form written $\oxmusafetyinnershape$,
which says that $\oxplaceexpr$ is $\oxmuta$-safe under $\oxtvarctx$ and
$\oxvarctx$, with \emph{reborrow exclusion list} $\overline{\oxplace}$, and
produces the loans
$\oxset{\overline{\oxloanpkg{\oxmuta}{\oxplaceexpr^\prime}}}$. Intuitively, we
use this reborrow exclusion list $\overline{\oxplace}$ to rule out precisely the
loan conflicts that arise from reborrowing as a programming pattern --- namely,
a reference conficting with loans from a reference it was reborrowed from.
Reborrowing is also what causes the borrow chain to contain more than just
$\oxplaceexpr$ for two reasons: (1)~we may be reborrowing from a reference whose
region has already lost precision, i.e.\ contains multiple loans, and thus we
cannot have perfect information about the reborrowed reference either, and
(2)~the reborrowed borrow chain will include an additional loan that records
that it was reborrowed (e.g. a unique loan for a reborrow from \rusti{x} would
appear as $\oxloanpkg{\oxmut}{\oxderef{\! x}}$). This second point is precisely
what produces the \emph{chain} aspect of the borrow chain since we can use this
information in each loan set to follow a series of consecutive reborrows.

Formally, the first rule, \oxname{O-SafePlace}, checks if a place $\oxplace$ is
$\oxmuta$-safe by looking at each loan in every region $\oxcprov^\prime$ in
$\oxvarctx$ and either (1) making sure that if either that loan or $\oxmuta$ is
$\oxmut$ then $\oxplace$ does not overlap with the loan; or (2) checking that
all references in $\oxvarctx$ with region $\oxcprov^{\prime}$ are in the
reborrow exclusion list (meaning we need not check if there is overlap with
$\oxplace$).

The next two rules check if a place expression $\oxplaceexpr$ is $\oxmuta$-safe,
decomposing the place expression into a place expression context
$\oxplaceexprctx$ (see \Fig{fig:term-syntax}) with $\oxderef{\oxplace}$ in the
hole. The last two lines of premises for both essentially ensure that either (1)
or (2) holds, but each one adds to the incoming reborrow exclusion list when
checking (2) by also including the place itself $\oxplace$ along with all of the
places $\oxplace_j$ that $\oxplace$ was borrowed from according to the loan set
$\oxcprov$ associated with its type. Both rules also check
$\oxmuta \oxsubtype \oxmuta_{\oxplace}$ (defined as the reflexive closure of
$\oximm \oxsubtype \oxmut$) in order to ensure that the reference has sufficient
permission to be used, preventing a dereference of a $\oxmut$ reference in a
$\oximm$ context.

Unlike \oxname{O-DerefAbs}, \oxname{O-Deref} is dereferencing a reference
$\oxplace$ with a concrete region $\oxcprov$. As such, we can look at the loans
present for $\oxcprov$ in the stack typing. These loans consist of both direct
loans to places $\oxplace_j$ which correspond to a possible origin for the
reference, and indirect loans to place expressions $\oxplaceexpr_j$ which
capture how this reference was reborrowed from other references. As such, when
we recursively check for the safety of these regions, we append the reborrow
origins (the $\oxplace_j$ prefixes of these $\oxplaceexpr_j$) to the reborrow
exclusion list. This means that they will not be considered as possible
conflicts in the rest of ownership safety. At the end, we union together the
borrow chains from all the possible origins to determine our final borrow chain.
We also include an additional loan
$\oxloanpkg{\oxmuta}{ \oxplaceexprctx[\oxderef{\oxplace}] }$ to indicate that
this use was reborrowed from $\oxderef{\oxplace}$.

\subsection{Typechecking \lang Programs}
\label{sec:typechecking}

\begin{figure*}
  \figuresize{ \vspace{-0.4em}
    \begin{flushleft}\oxtypjudgeform\end{flushleft}
    \begin{mathpar}
      \TuThreeTwo \and \TMove \and \TCopy \and \TBorrow \and \TLetProv \and
      \TBranch \qquad \TSeq \and \TLet \and \TDrop \and \TAssignDeref \and
      \TAssign \and \TClosure \and \TTuple \and \TWhile \and \TAbort
    \end{mathpar}
  }

  \caption{Selected \lang Typing Rules}
  \label{fig:typing}
\end{figure*}

\Fig{fig:typing} presents a selection of \lang typing rules. In every rule, we
highlight the expression being typechecked with a $\fbox{framebox}$. The shape
of our typing judgement is $\oxtypjudgeshape$: we typecheck $\oxexpr$ in a
global environment $\oxglobalctx$, type environment $\oxtvarctx$, temporary
typing $\oxkontctx$, and stack typing $\oxvarctx$, producing an updated stack
typing $\oxvarctx^\prime$ for typing the continuation of $\oxexpr$. These rules
rely on the region rewriting and outlives judgments (\Fig{fig:unification}),
which we'll discuss below, and the ownership safety judgment (\Fig{fig:safety}),
which we discussed in \Sec{sec:regions}. We elide the various well-formedness
judgments (for types, stack typings, etc.); \ifanonymous{see supplementary
  material.}{ see the appendix (\Sec{sec:well-formedness}). }

To best understand our typing judgment as a whole, it is useful to first know a
bit about what lies ahead in the metatheory (\Sec{sec:metatheory}). In our type
preservation proof, we need to maintain the well-typedness of values stored on
the stack in \lang. Since our values include closures which themselves may
introduce more aliasing, we then need to maintain our ownership safety judgment.
To make this possible, there are a number of restrictions that arise throughout
the type system on how regions annotate the program (and discussed further in
\Sec{sec:discussion}). We urge readers to keep this in mind.

\paragraph{Moving} In Oxide, as in Rust, owned values are moved or copied out of
a place $\oxplace$ when used, just as we saw in our first example in
\Sec{sec:ownership} where \rusti{pt} was moved to \rusti{x} and thus could not
be bound again to \rusti{y}. In order to move $\oxplace$, three conditions must
hold: (1)~$\oxplace$ must be able to be used $\oxmut$-ly (checked using the
ownership safety judgment in \Fig{fig:safety} from \Sec{sec:regions});
(2)~$\oxplace$ must have a sized and initialized (not dead) type $\oxsitype$ in
$\oxvarctx$; and (3)~this type $\oxsitype$ must be
$\oxkey{noncopyable}$.\footnote{We've elided definitions of $\oxkey{copyable}$
  and $\oxkey{noncopyable}$, but they're straightforward. Intuitively, a type is
  safe to copy if none of its constituent parts are unique. Thus, all types that
  don't contain a unique reference are copyable. Generic types are always
  non-copyable. In Rust, \oxkey{copyable} is actually the \rusti{Copy} trait,
  but \oxkey{copyable} can be thought of as special casing it.} If these
requirements hold, then we use a dagger to mark the place $\oxplace$ dead in the
continuation stack typing, preventing further expressions from reusing it.
Requirement~(1) is needed to ensure that we do not invalidate any existing
references to $\oxplace$ by moving it and requirement~(2) ensures that there is
currently data owned by $\oxplace$. If requirement~(3) does not hold, we'll copy
rather than move which permits more programs to typecheck since \oxname{T-Copy}
does not mark the copied place dead in the continuation stack typing as
\oxname{T-Move} does. Further, unlike moves (which are disallowed through
dereferences, leading to the restriction to places $\oxplace$ rather than place
expressions $\oxplaceexpr$ in \oxname{T-Move}), the copying variant
\oxname{T-Copy} can happen through a dereference and thus handles place
expressions generally.

\paragraph{Borrowing} As with moving, borrowing
$\oxref{\oxcprov}{\oxmuta}{\oxplaceexpr}$ relies on our ownership safety
judgment (\Fig{fig:safety}) with the $\oxmut$ and $\oximm$ modalities
corresponding precisely to the invariants of unique and shared pointers. Namely,
when $\oxmuta$ is $\oxmut$, we require that the place expression $\oxplaceexpr$
have no extant loans in $\oxvarctx$ and when $\oxmuta$ is $\oximm$, we require
no extant \emph{unique} loans. To actually know the type of the reference as a
whole, we also have to know the type of the place expression itself and we rely
on an auxillary judgment
$\oxcomputetynoprov{\oxtvarctx}{\oxvarctx}{\oxmuta}{\oxplaceexpr}{\oxxitype}$
(defined in the appendix) to compute the type $\oxxitype$ by starting with the
type of its root identifier and following the sequence of projections and
dereferences from $\oxplaceexpr$ through the type. For example, if we had the
place expression \rusti{*(x.0)} where x is a product of references with the type
\rusti{&'a uniq u32}, this judgment would produce the type \rusti{u32}. Much
like how \oxname{T-Move} updates the continuation stack typing to prevent
further uses of the moved place, \oxname{T-Borrow} updates the continuation
stack typing to associate the region $\oxcprov$ used for the borrow with the
loans that represent where the new reference may point (i.e.\ its provenance).
In many simple cases (such as borrowing a binding \rusti{x} with type
\rusti{u32}), there will be only one loan (corresponding to precise knowledge of
where it points), but as we will see with branching, these loan sets may grow
larger to account for information loss inherent to static analysis of reference
provenance.

\paragraph{Region Rewriting and Outlives} We examine region rewriting next
(\Fig{fig:unification}) since some of the typing rules discussed below require
it. The region rewriting judgment
$\oxtunify{\oxtvarctx}{\oxkontctx}{\oxvarctx}{\oxtype_1}{\oxtype_2}{\oxvarctx^\prime}$
says under $\oxtvarctx$, $\oxkontctx$, and $\oxvarctx$, we can rewrite the
regions from $\oxtype_1$ into the corresponding regions in $\oxtype_2$ which
will then be interpreted under $\oxvarctx^\prime$. We produce an output
$\oxvarctx^\prime$ with updated regions to be used when typing the continuation
after an appeal to region rewriting. The need for this judgment arises from the
need for values to be given the same type, in spite of the fact that they can be
annotated differently originally. Consider, for example, a branching term such
as \rusti{if cond { &'a uniq x } else { &'b uniq y }}. In this case, the type of
each side of the branch will be something like \rusti{&'a uniq u32} and
\rusti{&'b uniq u32} respectively, but we need to give an overall type for the
term. To deal with this, we pick one of the two types and use rewriting to write
the other one into the chosen type. For safety reasons, we pick the region with
the \emph{shortest} scope.

The rewriting judgment itself is fairly straightforward: it is reflexive and
transitive. Each rule for larger types recursively rewrites in any smaller
types, threading the output environment much like our typing judgment. The
actual rewriting portion arises, as one might expect, in the rule for references
(\oxname{RR-Reference}) which appeals to a judgment on regions which says that
the region from $\oxtype_1$ \emph{outlives} the region from $\oxtype_2$ while
performing the work required to enable the rewriting.

The outlives judgment (\Fig{fig:unification})
$\oxrunify{\oxtvarctx}{\oxvarctx}{\oxprov_1}{\oxprov_2}{\oxvarctx^\prime}$ says
under $\oxtvarctx$ and $\oxvarctx$, $\oxprov_1$ outlives $\oxprov_2$ rewriting
the latter in $\oxvarctx^\prime$ according to the mode $\oxrewritemode$. Every
region outlives itself (reflexivity). An abstract region outlives another if
there's a corresponding outlives relation in $\oxtvarctx$
(\oxname{OL-BothAbstract}) or if we can transitively put together outlives
relations from $\oxtvarctx$ (\oxname{OL-Trans}).\footnote{We do not need
  transitivity for concrete regions beyond what we can already conclude from the
  remaining \oxname{OL} rules.} \oxname{OL-CombineConcrete} (in $\oxcombine$
mode) says that $\oxcprov_1$ outlives $\oxcprov_2$ if it occurs earlier than
$\oxcprov_2$ in $\oxvarctx$. It also requires that there not exist any
references with either region which have been reborrowed
($\oxnotreborrowed{\oxvarctx}{\oxcprov_1}$ and
$\oxnotreborrowed{\oxvarctx}{\oxcprov_2}$, where \oxkey{rnrb} is an abbreviation
for region-not-reborrowed). This invariant ensures that value typing is
preserved under region rewriting. When this is the case, the output typing
$\oxvarctx^\prime$ is updated to associate $\oxcprov_2$ with the union of the
loans from both $\oxcprov_1$ and $\oxcprov_2$. The variant
\oxname{OL-CheckConcrete} (in $\oxnoop$ mode) has the same obligations, but
instead leaves the output unchanged. The behavioral difference between these
rules is the reason for the rewriting modes.

The last two rules say when a concrete region outlives an abstract one and vice
versa. In essence, a concrete region $\oxcprov$ can only outlive an abstract
region $\oxabsprov$ (\oxname{OL-ConcreteAbstract}) if $\oxcprov$ was
\emph{reborrowed}. The first two premises check for reborrowing: $\oxcprov$'s
loan set must be non-empty (otherwise there is no reborrow), and must consist
solely of place expressions $\overline{\oxplaceexpr}$ (since place expressions,
unlike places, contain dereferences, which identifies this as a reborrow instead
of a borrow). The third premise collects all the regions $\overline{\oxprov_i}$
that annotate any references dereferenced in each place expression
$\oxplaceexpr_i$ (see the type-computation judgment $\oxcomputetyshape$ in the
\ifanonymous{supplementary material}{appendix
  (\Sec{sec:additional-judgments})}), while the last premise ensures that all of
these outlive $\oxabsprov$. The final rule, \oxname{OL-AbstractConcrete}, says
that an abstract region \emph{always} outlives a concrete region. This is subtle
but makes sense because any abstract region $\oxabsprov$ is bound in a top-level
function (recall that closures don't abstract over regions), while a concrete
region $\oxcprov$ must be bound by $\oxkey{letrgn}$s \emph{inside} the function
body. Ultimately, any concrete region $\oxcprov^\prime$ that gets substituted
for $\oxabsprov$ upon application will already exist before $\oxcprov$ (even for
recursive calls), meaning it outlives $\oxcprov$.

\begin{figure*}
  \figuresize
  \begin{bnf}
    \oxrewritemodegrammar
  \end{bnf}

%\vspace{-3ex}

  \begin{flushleft}\oxtunifyform\end{flushleft}

  \begin{mathpar}
    \URefl \and \UTrans \and \UUninit \and \UTuple \and \USharedRef
  \end{mathpar}

  \vspace{1ex}
  \begin{flushleft}\oxrunifyform\end{flushleft}
  \vspace{-1ex}
  \begin{mathpar}
    \UReflProv \and \UTransProv \and \UAbsProvs \and \UCombineLocalProvs \and
    \UCombineLocalProvsUnrest \and \UCheckLocalProvs \and \OLocalAbsProvs \and
    \OAbsLocalProvs
  \end{mathpar}

  \vspace{-1em}
  \caption{Region Rewriting and Outlives Relations in \lang}
  \label{fig:unification}
  \vspace{-1em}
\end{figure*}

\paragraph{Branching and Sequencing} The next two rules illustrate how stack
typings are threaded through larger programs since the form of our typing
judgment requires each rule to specify its continuation's stack typing.
\oxname{T-Branch} uses the stack typing $\oxvarctx_1$ that we get from typing
the conditional $\oxexpr_1$ when typing each of the two branches. The type
$\oxsitype$ ascribed to the overall expression must be a supertype of the types
$\oxsitype_2$ and $\oxsitype_3$ of each branch and equal to one of them.
Additionally, branching uses a union operation $\oxintersect$ to combine the
output stack typings from each branch to produce the final stack typing
$\oxvarctx^{\prime}$ for the overall expression. $\oxintersect$ requires that
types of bound variables in the two stack typings be equal (which potentially
demands use of \oxname{T-Drop} when typing the branches), and unions the loan
sets for each region $\oxcprov$ from both stack typings (full definition in
technical appendix). Note that we only need to union stack typings with
identical domains --- we typecheck both branches under $\oxvarctx_1$ so they
produce output stack typings with the same domains (since $\oxkey{let}$ and
$\oxkey{letrgn}$ are the only means for introducing variables and regions, but
both are lexically-scoped), and region rewriting does not change the domain of
stack typings between its input and its output.

When typing $\oxseq{\oxexpr_1}{\oxexpr_2}$, we typecheck $\oxexpr_2$ under the
stack typing $\oxvarctx_1$ we got from typechecking $\oxexpr_1$. But,
importantly, we apply a metafunction $\oxgcloans{\oxkontctx}{\cdot}$ to
$\oxvarctx_1$ to empty out the loan sets of regions not used in $\oxvarctx_1$
before typing $\oxexpr_2$ because $\oxexpr_1$ may have been a unique reference
that is thrown away at runtime before moving on to $\oxexpr_2$. Without
\emph{garbage collecting} loans, \lang would reject programs that are safe and
accepted by Rust. Namely, this clearing allows us to handle the sort of ``early
dropping of references'' inherent to non-lexical lifetimes. Specifically,
$\oxgcloans{\oxkontctx}{\oxvarctx}$ empties out the loan set of each $\oxcprov$
that does not appear in any of the types in $\oxvarctx$ or $\oxkontctx$. The
full formal definition is present in the technical appendix.

\paragraph{Binding} In \lang, \oxname{T-Let} is interesting in three ways.
First, we allow rewriting in \oxname{T-Let} to a specified annotated type. This
rewriting allows us to change the regions in the computed type to match the
annotated type by conservatively combining the loans associated with each region
into the output, as described earlier in the section on region rewriting. Then,
similar to sequencing, \oxname{T-Let} uses the metafunction
$\oxgcloans{\oxkontctx}{\cdot}$ to eliminate any loans that might be unnecessary
as a result of $\oxexpr_1$ potentially being promoted to the annotated type
$\oxsitype_a$. Additionally, in the output stack typing from $\oxexpr_2$, we see
that our binding for $\oxid$ must have a dead type $\oxsdtype$ with the whole
binding being dropped in the overall stack typing $\oxvarctx_2$ output from
\oxname{T-Let} (since the scope of $\oxid$ ends at that point). The requirement
that the type be dead means we must have either used \oxname{T-Move} to move out
of that binding or we must have explicitly used \oxname{T-Drop} on $\oxid$ in
the derivation for $\oxexpr_2$, and can be thought of as a formalization of the
``resource acquisition is initialization'' pattern~\cite{stroustrup94:raii}
since we are explicitly requiring a first-in-last-out allocation/deallocation
pattern and require everything to have been used in either a move or a drop
before we can end its scope.

\paragraph{Drop} As alluded to in the previous two paragraphs, \lang has a rule
called \oxname{T-Drop} which is used non-deterministically during typechecking
to mark a particular place $\oxplace$ as being dead. This rule corresponds
roughly to a conventional weakening rule where $\oxplace$ ``doesn't exist'' (in
this case, is dead) in the premise, but exists in the conclusion. The main
difference is that while the data is thought to be deallocated, the name is
still in-scope to be dropped when its scope ends in \oxname{T-Let}.

\paragraph{Assignment} Assignment is interesting in a few ways. First,
assignment is broken up into two rules \oxname{T-Assign} and
\oxname{T-AssignDeref} where the former is able to assign to a place $\oxplace$
that is dead, and the latter is able to assign to a place through a reference
(i.e. by using dereferencing). The basic structure of each rule is the same. For
both rules, we typecheck the new expression to be assigned, look up the type of
the place we're assigning to (a lookup in \oxname{T-Assign} and a type
computation in \oxname{T-AssignDeref}), check compatibility of the new
expression's type with the type of the place we're assigning to, and then
finally check that it's safe to use the place we're assigning to uniquely
according to ownership safety.

The differences between the two play a fundamental role in allowing us to
appropriately model how assignment works in Rust. Notably, the region rewriting
judgment in \oxname{T-Assign} uses the checking mode (denoted $\oxnoop$) to
limit how conservative the borrow checker need be after an assignment. As
discussed earlier, this mode does not change its output environment (thus,
$\oxvarctx^\prime = \oxrsub{\oxvarctx_1}{\oxderef{\oxplace}}$ in
\oxname{T-Assign}). This is okay in context because after the typing rule is
done, the type of $\oxplace$ is updated to $\oxsitype$ (the type of its new
value) in the continuation. A similar update in \oxname{T-AssignDeref} would
entail updating the types of arbitrarily many bindings, and so instead the more
conservative combine mode (denoted $\oxcombine$) is used. Further, in
\oxname{T-Assign}, we employ the operation
$\oxrsub{\oxvarctx_1}{\oxderef{\oxplace}}$ for the input environment to the
region rewriting judgment. This operation is defined to remove any loans
prefixed by $\oxderef{\oxplace}$ from every loan set in $\oxvarctx_1$. The
$\rsub$ operation (informally called the ``kill rules'' in Rust) amounts to
erasing reborrowing relationships that no longer hold as a result of this
assignment.

Concretely, consider an environment with two references, one named \rusti{x}
with type \rusti{&'x uniq u32} and another named \rusti{y} reborrowed from
\rusti{x} with type \rusti{&'y uniq u32}. This means in our environment we would
have loan sets that look something like
$\oxloanctxentry{\rusti{'x}}{\oxset{\overline{\oxloan}}}$ and
$\oxloanctxentry{\rusti{'y}}{ \oxset{\overline{\oxloan} \oxcomma \oxloanpkg{\oxmut}{\oxderef{\!\oxid}}} }$.
If we were to then assign to \rusti{x} some unrelated reference, the kill rules
would delete the loan $\oxderef{\oxid}$ in the loan set of \rusti{'y} since
after the assignment ran, the two would represent distinct and disjoint
references.

The last difference between the two assignment rules is presence of an
additional obligation in \oxname{T-Assign}:
$\oxrgnuniqueto{\oxsxtype}{\oxplace}{\oxvarctx_1}$. This obligation means that
there are not any places in $\oxvarctx_1$ that share the outermost region of
$\oxsxtype$ (i.e. if $\oxsxtype = \oxtref{\oxcprov}{\oxmuta}{\oxxitype}$, then
$\oxcprov$ would be its outermost region). This allows us to guarantee that the
garbage collection discussed for the sequencing rule will always clear out this
outermost region $\oxcprov$ before the subsequent expression is typechecked,
helping us greatly in our proofs. It's important to note that this obligation
pertains to annotations added to the \lang program compared to Rust, and so it
only limits the patterns of region annotation that can be applied, rather than
the space of Rust programs that can be typechecked in \lang.

\begin{figure*}
  \figuresize
  \begin{mathpar}
    \TAppClosure
  \end{mathpar}

  \vspace{-1em}
  \caption{\lang Typing Rule for Application}
  \label{fig:typing-app}
  \vspace{-1em}
\end{figure*}

\paragraph{Closures and Application}
Closures in \lang correspond to \emph{move closures} in Rust which move or copy
their free variables from the outer environment into the
closure.\footnote{Rust's standard closures implicitly introduce borrowed
  temporaries for all the free variables. We can recover this behavior via a
  simple, local transformation to move closures.} As such, \oxname{T-Closure}
must compute the captured frame by looking at the free variables (and free
regions) of the closure's body, and it must mark dead (add daggers to the types
of) any variables in the stack typing with \emph{non-copyable} types. The
captured frame is suspended over the arrow in the function type to keep track of
the fact that the data caught up in the closure is still alive (and thus must be
considered in ownership safety). We elide the rule for top-level function
definitions, which gives a function $\oxfnname$ the type that $\oxfnname$ is
annotated with in $\oxglobalctx$, relying on well-formedness of $\oxglobalctx$
to know that this is okay.

The rule for application (\oxname{T-AppClosure} in \Fig{fig:typing-app}) is
roughly as one would expect: we typecheck the function, and then the arguments,
threading through the environments. However, at each step, we do a region
rewriting using the unrestricted combine mode (denoted $\oxcombineunrest$) for
the computed argument types to the annotated parameter types. This unrestricted
mode allows us to push loans from the surrounding context into the regions used
in the closure's signature, a behavior ruled out by the conventional combine
mode ($\oxcombine$) because they would enable degenerate annotation patterns
that make it difficult to prove soundness. For top-level functions, we have an
additional rule \oxname{T-AppFunction} in the technical appendix that does not
use a rewriting for the types, but (1)~substitutes all frame, type, and region
variables in the types and (2)~checks that any outlives bounds specified on the
function signature hold (via outlives in \Fig{fig:unification}). This does not
appear in \oxname{T-AppClosure} since closures cannot be polymorphic, nor can
they possess where bounds.

\paragraph{Values and Aggregates} The typing rules for base types
(\oxname{T-u32}, \oxname{T-True}, \oxname{T-False}, etc.) are standard, and
leave the type environment unchanged in their output. Aggregate structures like
tuples check the types of their components while threading through the
environments in left-to-right order. This left-to-right ordering for
typechecking corresponds to the ordering implemented by Rust's typechecker and
borrow checker. One subtlety to note (discussed already with the last example in
\Sec{sec:envs}) is that when typechecking a component $e_i$ of the tuple, we add
the types of all earlier tuple components to the temporary typing $\oxkontctx$.
This is needed because $e_i$ might be a let-binding or sequencing expression
that invokes $\oxgcloans{\oxkontctx}{\cdot}$ on its environment during
typechecking. If the types of the earlier tuple components we've just
typechecked aren't present in the environment, to serve as roots for the
regions mentioned in those types, we may end up incorrectly garbage collecting
the loans that these regions map to. This would make programs with subsequent
borrows typecheck even when the borrow should not be allowed. The elided typing
rule for arrays is similar.

The formalism of \lang omits a specific treatment of structs, but we note that
they are essentially the same as tuples, only featuring a tag that must also be
checked. Our implementation which we discuss in \Sec{sec:tested-semantics}
relies on exactly this approach to support structs.

\paragraph{Remaining Rules} The remaining rules in \Fig{fig:typing} are
straightforward or covered earlier. Elided typing rules all concern arrays and
are given in the \ifanonymous{supplementary material}{ technical appendix
  (\Sec{sec:typing}) }.

\subsection{Operational Semantics}
\label{sec:dynamics}

\begin{figure}
  \figuresize
  \begin{bnf}
    \oxreferentgrammar \\[0em]
    \oxptrexprgrammar \\[0em]
    \oxevalctxgrammar \\[1em]
    \oxvaluegrammar \oxptrextension \\[0em]
    \oxvaluectxgrammar \\[0em]
    \oxstoregrammar \\[0em]
    \oxstackframegrammar \\
  \end{bnf}

  \vspace{-1em}
  \caption{\lang Syntax Extensions for Dynamics}
  \label{fig:dynamics-syntax}
  \vspace{-1em}
\end{figure}

\begin{figure*}
  \figuresize
  \begin{flushleft}\oxnormform \hspace{2em}
    % read: ``$\oxplaceexpr$ computes to $\oxreferent$, which maps to
    % $\oxvalue$ in $\oxstore$.'' \\[1em]
    $\oxnorm{\oxstore}{\oxplaceexprctx[\oxid]}
    {\oxreferent}{\oxvaluectx}{\oxvalue} ~~\defeq~~
    \oxnorminner{\oxstore}{\oxplaceexprctx}{\oxid}
    {\oxreferent}{\oxvaluectx}{\oxvalue}$.
  \end{flushleft}

  \begin{flushleft}\oxnorminnerform \hspace{1em}
    read: ``$\oxreferent$ in a context $\oxplaceexprctx$ computes to
    $\oxreferent^\prime$ which maps to $\oxvalue$ in
    $\oxstore$.''\end{flushleft}
  \begin{mathpar}
    \PId \qquad \PProj \qquad \PDerefPtr
  \end{mathpar}

  \vspace{0.5em}

  \begin{flushleft}\oxreduceform\end{flushleft}

  \begin{mathpar}
    \EMove \and \ECopy \and \EBorrow \and \ESeq \quad \ELetProv \quad \EAssign
    \and \ELet \and \EShift \and \EClosure \and \EApp \and \EFramed \and \EWhile
  \end{mathpar}

  \vspace{-1em}
  \caption{Selected Place Expression Evaluation Rules (top) and Reduction Rules
    (bottom)}
  \label{fig:dynamics}
  \vspace{-1em}
\end{figure*}

For our operational semantics, we extend the syntax of \lang in
\Fig{fig:dynamics-syntax} with terms that only arise at runtime. First, to be
able to specify what ``address'' a pointer points to, we introduce an abstract
form of memory addresses called \emph{referents}. Referents $\oxreferent$
essentially record what the offsets are from a variable on the stack in order to
specify a precise ``memory address,'' (e.g., a particular element of an array or
tuple, or a particular slice of an array). We also include some administrative
forms: (1)~$\oxframed{\oxexpr}$ and $\oxshift{\oxexpr}$ which are used when
evaluating application and let bindings discussed below, (2)~$\oxuninit$ (the
dead value), and (3)~$\oxsliceval{\oxvalue_1 \oxdotsc \oxvalue_n}$ which is a
dynamically-sized slice of an array. Then, we introduce value forms including
pointers to referents, and closures packaged with their environment
$\oxstackframe$. \Fig{fig:dynamics-syntax} also includes stacks $\oxstore$ as
a sequence of stack frames $\oxstackframe$, and value contexts $\oxvaluectx$
which allow array values to be decomposed with multiple holes when dealing with
slices.

In \Fig{fig:dynamics}, we present a selection of our small-step operational
semantics which is defined using Felleisen and Hieb-style left-to-right
evaluation contexts~\cite{felleisen92:eval-contexts} over configurations of the
form $\oxconfigurationshape$. Since our semantics uses referents $\oxreferent$
as an abstract version of memory addresses, some of our rules rely on a notion
of place-expression evaluation,
$\oxnorm{\oxstore}{\oxplaceexpr}{\oxreferent}{\oxvaluectx}{\oxvalue}$
(\Fig{fig:dynamics}, top), which should be read as: $\oxplaceexpr$ evaluates to
$\oxreferent$, which maps to $\oxvalue$ with a surrounding context $\oxvaluectx$
in $\oxstore$.

The evaluation rules are straightforward: \oxname{E-Move} returns a value by
moving it off of the stack $\oxstore$, replacing it with $\oxuninit$.
\oxname{E-Copy} copies the value from the stack. \oxname{E-Borrow} creates a
pointer value to the referent $\oxreferent$. Branching is completely standard,
hence elided. Assignment, similar to \oxname{E-Copy} and \oxname{E-Borrow}, uses
the place-expression evaluation rules, but instead cares specifically about the
value context $\oxvaluectx$, rather than the value $\oxvalue$. Assignment also
decomposes the computed referent $\oxreferent$ to get its root identifier
$\oxid$. Then, it updates the stack by maintaining this value context when it
updates $\oxid$ (mapping it to $\oxvaluectx[\oxvalue]$).

\paragraph{Binding and the Stack} Bindings are interesting in that they
introduce our two administrative forms, $\oxframed{\oxexpr}$ and
$\oxshift{\oxexpr}$. For instance, in \oxname{E-Let}, we step to
$\oxshift{\oxexpr}$ rather than $\oxexpr$ alone in order to ensure that the
binding for $\oxid$ is well-scoped and ends when it should (seen in
\oxname{E-Shift}). In \oxname{E-AppClosure}, we similarly step to
$\oxframed{\oxexpr}$ to ensure that after evaluating the body of the closure we
drop the stack frame from that function call (seen in \oxname{E-Framed}). Both
\oxname{E-Shift} and \oxname{E-Framed} rely crucially on the fact that our stack
$\oxstore$ is ordered --- they must match the most recent entry.

\subsection{Well-typed \lang programs won't go wrong!}
\label{sec:metatheory}

We prove syntactic type safety for \lang using progress and
preservation~\cite{wright92:progress-preservation}.

\begin{oxlemma}{Progress}{lemmap:progress}
  If \oxtypjudge{\oxglobalctx}{\oxemptyctx}{\oxkontctx}{\oxvarctx}{\oxexpr}
  {\oxsitype}{\oxvarctx^\prime} and
  \oxstorevalidity{\oxglobalctx}{\oxvarctx}{\oxstore}, then either $\oxexpr$ is
  a value, $\oxexpr$ is an \oxabort{\,\dots\,}, or
  $\oxexists \oxstore^\prime, \oxexpr^\prime. \;
  \oxreduce{\oxglobalctx}{\oxstore}{\oxexpr}{\oxstore^\prime}{\oxexpr^\prime}$.
\end{oxlemma}

The Progress lemma says that if we can typecheck $\oxexpr$ under a valid global
environment $\oxglobalctx$, temporary typing $\oxkontctx$, and stack typing
$\oxvarctx$, \emph{and} we have a stack $\oxstore$ that satisfies this stack
typing $\oxvarctx$, then either $\oxexpr$ is a value, an \rusti{abort!}
expression, or we can take a step. Our stack typing judgment
$\oxstorevalidity{\oxglobalctx}{\oxvarctx}{\oxstore}$ says that each value in
the stack $\oxstore$ has the corresponding type attributed to it in the typing
$\oxvarctx$. The proof proceeds by induction on the typing derivation for
$\oxexpr$, and relies on a Canonical Forms lemma and
\Lemma{lemmap:place-exprs-reduce} which says that place expressions can be
reduced at runtime to values with their computed types. This lets us apply the
rules for moves, copies, borrowing, and assignment.

\begin{oxlemma}{Place Expressions Reduce}{lemmap:place-exprs-reduce}
  If $\oxcomputetynoprov{\oxtvarctx}{\oxvarctx}{\oxmuta}
  {\oxplaceexpr}{\oxxitype}$ and
  $\oxstorevalidity{\oxglobalctx}{\oxvarctx}{\oxstore}$, then
  $\oxnorm{\oxstore}{\oxplaceexpr}{\oxreferent}{\oxvaluectx}{\oxvalue}$ and
  $\oxtypjudge{\oxglobalctx}{\oxtvarctx}{\oxkontctx}{\oxvarctx}
  {\oxvalue}{\oxxitype}{\oxvarctx}$.
\end{oxlemma}

Our proof of \Lemma{lemmap:place-exprs-reduce} (ifanonymous{Lemma 5.3}{\Lemma{lemma:place-exprs-reduce}} in appendix) relies on
the shared inductive structure of type computation
$\oxcomputetynoprov{\oxtvarctx}{\oxvarctx}{\oxmuta}{\oxplaceexpr}{\oxxitype}$
and place expression evaluation
$\oxnorm{\oxstore}{\oxplaceexpr}{\oxreferent}{\oxvaluectx}{\oxvalue}$.

\begin{oxlemma}{Preservation}{lemmap:preservation}
  If \oxtypjudge{\oxglobalctx}{\oxemptyctx}{\oxkontctx}{\oxvarctx}{\oxexpr}
  {\oxsitype_1}{\oxvarctx_f} and
  \oxstorevalidity{\oxglobalctx}{\oxvarctx}{\oxstore} and
  $\oxwfkontctx{\oxglobalctx}{\oxvarctx}{\overline{\oxvalue}}{\oxkontctx}$ and
  \oxreduce{\oxglobalctx}{\oxstore}{\oxexpr}{\oxstore^\prime}{\oxexpr^\prime}, then there
  exists $\oxvarctx_i$ such that
  \oxstorevalidity{\oxglobalctx}{\oxvarctx_i}{\oxstore^\prime} and
  $\oxwfkontctx{\oxglobalctx}{\oxvarctx_i}{\overline{\oxvalue}}{\oxkontctx}$ and
  \oxtypjudge{\oxglobalctx}{\oxemptyctx}{\oxkontctx}{\oxvarctx_i}{\oxexpr^\prime}
  {\oxsitype_2}{\oxvarctx_f^\prime} and
  $\oxtunify[\oxcombine]{\oxemptyctx}{\oxvarctx_f^\prime}
  {\oxsitype_2}{\oxsitype_1}{\oxvarctx_s}$ and there exists $\oxvarctx_o$ such
  that $\oxvarctx_f = \oxvarctx_s \oxintersect \oxvarctx_o$.
\end{oxlemma}

The Preservation lemma says that if $\oxexpr$ has type $\oxsitype_1$ under a
valid global environment $\oxglobalctx$, temporary typing $\oxkontctx$, and
stack typing $\oxvarctx$, \emph{and} we have a stack $\oxstore$ that satisfies
$\oxvarctx$ and a sequence of values $\overline{\oxvalue}$ that satisfies the
temporary typing $\oxkontctx$, \emph{and} we know that $\oxexpr$ can take a step
under $\oxstore$ to the new configuration
$\oxconfiguration{\oxstore^\prime}{\oxexpr^\prime}$, then the following conditions hold.
Our updated stack $\oxstore^\prime$ satisfies $\oxvarctx_i$, our sequence of
temporary values $\overline{\oxvalue}$ continue to satisfy $\oxkontctx$, and our
new expression $\oxexpr^\prime$ typechecks with the type $\oxsitype_2$. In each of
these judgments, we use an intermediate stack typing $\oxvarctx_i$ that
corresponds to the changes the evaluated portion of code made to the
environment. Rather than constrain the type to be the same as the type of our
original $\oxexpr$, our Preservation lemma allows $\oxsitype_2$ to differ in its
regions by the region rewriting judgment, since evaluation potentially can lead
to a type having more precise regions. Further, the output stack typing from
typechecking $\oxexpr^\prime$ is threaded through the rewriting and then ultimately
said to union with some other stack typing $\oxvarctx_o$ in order to capture the
relationship between the old output environment $\oxvarctx_f$ and the new one
$\oxvarctx_f^\prime$ when we have taken a step into one side or the other of a
branch.

As discussed in \Sec{sec:typechecking}, the most challenging part of proving
preservation is in showing that the various changes to the environment preserve
the well-typedness of values, with closures in particular posing the greatest
issue. Indeed, we believe it's clear that a formalization of Rust without
closures misses a significant piece of the language's essence since closures
interact consistently with all parts of the formalism. To that end, the
proof of preservation uses several families of lemmas that
follow the same pattern: since our preservation theorem has to maintain the
well-typedness of the stack $\oxstore$ and temporary values
$\overline{\oxvalue}$, we must show that various judgments in our system
preserve the well-typedness of values. We will highlight these lemmas here,
noting that each require sublemmas for expressions in closure bodies
remaining well-typed which subsequently requires that region rewriting and
ownership safety judgments are preserved by these judgments.

\begin{oxlemma}{Values are Well-Typed after Region Rewriting}
  {lemmap:value-typing-region-rewriting}
  \ifanonymous{(Lemma 5.12}{(\Lemma{lemma:subtyping-value-typing}} in appendix) \hfill \\
  If \ $\oxtypjudge{\oxglobalctx}{\oxtvarctx}{\oxkontctx}{\oxvarctx}
  {\oxvalue}{\oxtype}{\oxvarctx}$ and
  $\oxtunify{\oxtvarctx}{\oxvarctx}{\oxtype_2}{\oxtype_1}{\oxvarctx^\prime}$
  then $\oxtypjudge{\oxglobalctx}{\oxtvarctx}{\oxkontctx}{\oxvarctx^\prime}
  {\oxvalue}{\oxtype}{\oxvarctx^\prime}$.
\end{oxlemma}

\begin{oxlemma}{Values are Well-Typed after Drop/GC-Loans}
  {lemmap:value-typing-related-envs}
  (\ifanonymous{Lemma 5.25}{\Lemma{lemma:value-typing-related-envs}}) \hfill \\
  If \ $\oxsubtypectx{\oxglobalctx}{\oxemptyctx}{\oxvarctx} {\oxvarctx^\prime}$ and
  $\oxtypjudge{\oxglobalctx}{\oxemptyctx}{\oxemptyctx}{\oxvarctx}
  {\oxvalue}{\oxvarctx(x)}{\oxvarctx}$, then
  $\oxtypjudge{\oxglobalctx}{\oxemptyctx}{\oxemptyctx}{\oxvarctx^\prime} {\oxvalue}{\oxvarctx^\prime(x)}{\oxvarctx^\prime}$.
\end{oxlemma}

\begin{oxlemma}{Values are Well-Typed under Well-Typed Extensions}
  {lemmap:value-typing-extension}
  (\ifanonymous{Lemma 5.34}{\Lemma{lemma:values-well-typed-extension}}) \hfill \\
  If \ $\oxtypevalidity{\oxglobalctx}{\oxemptyctx}{\oxvarctx}{\oxsitype_\oxid}$ and
  $\forall \oxcprov \in \oxfprovs{\oxsitype_\oxid}. \;
  \oxnotreborrowed{\oxvarctx}{\oxcprov}$ and
  $\oxtypjudge{\oxglobalctx}{\oxemptyctx}{\oxkontctx}{\oxvarctx}
  {\oxvalue}{\oxsitype}{\oxvarctx}$, then
  $\oxtypjudge{\oxglobalctx}{\oxemptyctx}{\oxkontctx}{ \oxextendctx{\oxvarctx}{
      \oxvarctxentry{\oxid}{\oxsitype_\oxid} } }{\oxvalue}{\oxsitype}{
    \oxextendctx{\oxvarctx}{ \oxvarctxentry{\oxid}{\oxsitype_\oxid} } }$.
\end{oxlemma}

\begin{oxlemma}{Values are Well-Typed after Assignment}
  {lemmap:values-well-typed-assignment}
  (\ifanonymous{Lemma 5.40}{\Lemma{lemma:values-well-typed-assignment}}) \hfill \\
  If \ $\oxlookup{\oxvarctx}{\oxplace_a}{\oxsxtype}$ $\wedge$
  $\oxtunify[\oxnoop]{\oxtvarctx}{\oxkontctx}{
    \oxrsub{\oxvarctx}{\oxderef{\oxplace_a}}
  }{\oxsitype}{\oxsxtype}{\oxvarctx^\prime}$ $\wedge$
  $\oxmusafety{\oxemptyctx}{\oxkontctx}{\oxvarctx^\prime}{\oxmut}{\oxplace_a}{
    \oxset{\oxloanpkg{\oxmut}{\oxplace_a}}
  }$ $\wedge$
  $\oxtypjudge{\oxglobalctx}{\oxemptyctx}{\oxkontctx}{\oxvarctx}
  {\oxvalue}{\oxtype}{\oxvarctx}$, then
  $\oxtypjudge{\oxglobalctx}{\oxemptyctx}{\oxkontctx}{
    \oxgcloans{\oxkontctx}{\oxtupdate{\oxvarctx^\prime}{\oxplace_a}{\oxsitype}}
  }{
    \oxvalue
  }{\oxtype}{
    \oxgcloans{\oxkontctx}{\oxtupdate{\oxvarctx^\prime}{\oxplace_a}{\oxsitype}}
  }$.
\end{oxlemma}

\begin{oxlemma}{Values are Well-Typed under Safe Loan Updates}
  {lemmap:value-typing-loan-update}
  (\ifanonymous{Lemma 5.62}{\Lemma{lemma:values-well-typed-loan-update}}) \hfill \\
  If \ $\oxmusafety{\oxemptyctx}{\oxkontctx}{\oxvarctx}{\oxmuta}{\oxplaceexpr}{
    \oxset{\overline{\oxloan}} }$ and
  $\oxnotinclosure{\oxkontctx}{\oxvarctx}{\oxcprov}$ and
  $\oxlookup{\oxvarctx}{\oxcprov}{\emptyset}$ and
  $\oxtypjudge{\oxglobalctx}{\oxemptyctx}{\oxkontctx}{\oxvarctx}
  {\oxvalue}{\oxsitype}{\oxvarctx}$, then
  $\oxtypjudge{\oxglobalctx}{\oxemptyctx}{\oxkontctx}{ \oxupdate{\oxvarctx}{
      \oxloanctxentry{\oxcprov}{ \oxset{\overline{\oxloan}} } }
  }{\oxvalue}{\oxsitype}{ \oxupdate{\oxvarctx}{ \oxloanctxentry{\oxcprov}{
        \oxset{\overline{\oxloan}} } } }$.
\end{oxlemma}

This pattern of a value typing lemma demanding a whole family of related lemmas
is consistent throughout the supporting lemmas, and they eventually rely on an
ownership safety lemma where we need to actually consider how the various
judgments each change the environment and argue that the separation provided by
closures --- including for regions --- is sufficient to prevent those changes
from breaking ownership safety derivations in the closure bodies themselves. In
each case, the hardest part of proving each family is ensuring that the
induction hypothesis for the ownership safety lemma was rich enough to address
the changes effects on the loans in the environment and the reborrow exclusion
list discussed in \Sec{sec:regions}.

With \Lemma{lemmap:progress} and \Lemma{lemmap:preservation} in hand, we also
prove a conventional type safety theorem as a corollary. However, we note that
Preservation itself represents a more interesting metatheoretic result because
it requires us to show that all of the type system invariants (namely, the
aliasing requirements) are maintained throughout the execution of well-typed
programs.

\begin{figure*}
  \figuresize
  \include{test-table}

  \vspace{-1em}
  \caption{Tested Semantics Results}
  \label{fig:test-results}
  \vspace{-1em}
\end{figure*}

\subsection{Tested Semantics}
\label{sec:tested-semantics}

We set out at the onset to solve a particular problem --- there is no high-level
specification of the Rust programming language and its borrowchecker. If there
were, this would be the point where we might present a proof that every
expression that typechecks in \lang also typechecks in Rust and vice versa.
Since doing that is not possible, we follow \citet{guha10:essence-js} in
developing a \emph{tested semantics} for \lang. We built an implementation of
our \lang typechecking algorithm, \typechecker, alongside a compiler, \compiler,
from a subset of Rust (with a small number of additional annotations) to \lang.
In addition to the features described in \Sec{sec:formalization}, our
implementation supports \rusti{struct}s by treating them as tagged tuples or
records. Together, \compiler and \typechecker allowed us to use tests from the
official borrow checker (\rusti{borrowck}) and non-lexical lifetime
(\rusti{nll}) test suites to validate \lang against Rust's implementation,
\rustc. The results of this testing are summarized in \Fig{fig:test-results}.

For the 208 passing tests, we can compile the test case into \lang with
\compiler and then use \typechecker to either successfully typecheck the
program or to produce a type error. We compare this typechecking result to the
expected behavior according to the \rustc test suite. All 208 tests either type
check when \rustc does so, or produce an error corresponding to the error
produced by \rustc.

The remaining 407 tests were taken out of consideration on the basis of being
out-of-scope for this work. There were 20 categories for exclusion, the majority
of which had fewer than 10 applicable tests. \Fig{fig:test-results} includes the
6 largest categories: (1)~heap allocation, (2)~out-of-scope libraries,
(3)~enumerations, (4)~statics and constants, (5)~traits, and (6)~uninitialized
variables. One specialized category (multithreading) was folded into
out-of-scope libraries in this table, with miscellaneous aggregating the
remaining smaller categories: control flow, casting, first-class constructors,
compiler internals dumping, function mutability, inline assembly, macros, slice
patterns, two-phase borrows, uninitialized variables, universal function-call
syntax, unsafe, and variable mutability.

Combined, heap allocation and out-of-scope libraries (of which the former is a
specialization of the latter) make up for the largest excluded category with 103
tests, and can be extended in future work using the strategy outlined by
\citet{weiss19:oxide}. The next largest category, traits, accounts for 93 tests.
Though the trait system is in some ways novel, the bulk of its design is rooted
in the work on Haskell typeclasses and their extensions. As such, we feel that
they are not an \emph{essential} part of Rust, though exploring the
particularities of their design may be a fruitful avenue for future work on
typeclasses. We are working on extending our implementation with sums to support
enumerations, and they are already present in the formalism. Many of the other
categories describe features (e.g., macros, control flow, casting, statics, and
constants) that are well-studied in the literature, and in which we believe Rust
has made relatively standard design choices.

The last issue to discuss involving the tested semantics is the aforementioned
annotation burden. This burden comes directly out of the syntactic differences
between \lang and Rust, and so are overall rather minor. The most immediately
apparent need is to provide a region annotation on borrow expressions, which we
handle using Rust's compiler annotation support. In our tests, a borrow
expression like $\oxref{\texttt{'a}}{\,\oxmut\,}{\texttt{x}}$ appears as
\rusti{#[lft="a"] &mut x}. However, we reduce the need for this by automatically
generating a fresh local region for borrow expressions without an annotation.
This suffices for the majority of expressions without change. Relatedly, one
might also expect to see the introduction of \oxkey{letrgn} throughout. To
alleviate the need for this, our implementation automatically binds free region
at the beginning of each function body.

The other main change we had to make relates to the use of explicit environment
polymorphism in \lang. In Rust, every closure has a unique type without a syntax
for writing it down. To work with higher-order functions, these closures
implement one of three language-defined traits (\rusti{Fn}, \rusti{FnMut}, and
\rusti{FnOnce}) which can be used as bounds in higher-order functions. We
compile the use of these trait bounds to environment polymorphism in a
straight-forward manner (turning instances of the same \rusti{Fn}-bound
polymorphic type into uses of function types with the same environment
variable), but need to introduce a way of writing down which environment to use
at instantiation. We use a compiler annotation (\rusti{#[envs(c1, ..., cn)]}) on
applications which says to instantiate the environment variables with the
captured environments of the types of these bindings. If the bindings are
unbound or not at a function type, we produce an error indicating as much.

Aside from these two changes, there are a handful of smaller changes that we
made by hand to simplify implementation of \compiler and \typechecker, though
the need for these could be obviated with more work. Our implementation does not
support method call syntax, and so we translate method definitions (which take
\rusti{self}, \rusti{&self}, or \rusti{&mut self} as their first argument) into
function definitions with a named first argument at the method receiver's type.
Relatedly, some of the tests used traits in a trivial way to define methods
polymorphic in their receiver type. Like other methods, we translated these into
function definitions, but used a polymorphic type for the receiver. \rustc also
allows for a number of convenient programming patterns (like borrowing from a
constant, e.g. \rusti{&0}) which are not supported by our implementation. To
deal with these cases, we manually introduced temporaries (a process that \rustc
does automatically). As a simplification for the typechecker, \typechecker only
reports the first error that occurs in the program. To ensure that we find a
correspondence between all errors, we split up test files with multiple errors
into one file per test.

Finally, an earlier version of our implementation required type annotations on
all let bindings, and so currently many tests include fully-annotated types. We
later realized our typing judgment is very-nearly a type \emph{synthesis}
judgment as in bidirectional typechecking, and changed the implementation to
support unannotated bindings by using the type synthesized for the expression
being bound. This works for all expressions except \oxkey{abort!} which can
produce any type and so requires annotation.

\ifanonymous{
  We provide a snapshot of the tooling as anonymized supplementary materials to
  our submission, including the entire \rusti{borrowck} and \rusti{nll} test
  suites and all disqualified tests categorized appropriately. The test suite
  features a full, anonymized \oxkey{git} history that captures the changes
  summarized above.
}{}

%%% Local Variables:
%%% mode: latex
%%% fill-column: 80
%%% End:
